\part{Das Werk als Ganzes}

\chapter{Programm und Werk}\label{kap:werk}

\begin{merksatz}
\boxtitel{Der Unterschied}
\begin{center}
\renewcommand{\arraystretch}{1.5}
\begin{tabular}{@{}P{2.6cm}P{2.6cm}P{3.0cm}@{}}
\toprule
& \textbf{Programm} & \textbf{Werk}\\
\midrule
Läuft & \textbf{einmal, irgendwo} & \textbf{dauernd, an vielen Orten}\\
Fehler & \textcolor{mittel}{Absturz} & \textcolor{akzent}{\textbf{Teilausfall}}\\
Ändern & neu übersetzen & \textcolor{akzent}{im laufenden Betrieb}\\
Zustand & im Speicher & \textcolor{akzent}{\textbf{verteilt, dauerhaft}}\\
Zuständigkeit & \textbf{einer} & \textcolor{akzent}{\textbf{mehrere}}\\
\bottomrule
\end{tabular}
\end{center}

\begin{center}
\begin{minipage}{0.88\textwidth}
\kursiv{Die zweite Zeile ist der wichtigste Unterschied: Ein Programm läuft oder es läuft nicht. Ein Werk kann zu vierzig Prozent laufen --- und diese Möglichkeit ist der Ursprung fast aller Schwierigkeiten, die dieser Band behandelt.}
\end{minipage}
\end{center}

\textbf{Was der Teilausfall bedeutet.}

\begin{center}
{\small
\renewcommand{\arraystretch}{1.45}
\begin{tabular}{@{}P{2.8cm}P{2.6cm}P{2.6cm}@{}}
\toprule
\textbf{Lage} & \textbf{Im Programm} & \textbf{Im Werk}\\
\midrule
Aufruf schlägt fehl & \textcolor{gruen}{Fehlerwert} & \textcolor{akzent}{\textbf{oder Stille}}\\
Empfänger weg & \textcolor{akzent}{gibt es nicht} & \textcolor{akzent}{Normalfall}\\
Antwort verspätet & \textcolor{akzent}{gibt es nicht} & \textcolor{akzent}{\textbf{nicht von \enquote{weg} zu trennen}}\\
Halb ausgeführt & \textcolor{gruen}{selten} & \textcolor{akzent}{\textbf{häufig}}\\
\bottomrule
\end{tabular}}
\end{center}

\begin{center}
\kursiv{Die dritte Zeile ist die grundlegende: Ein Aufruf, der nicht antwortet, kann verlorengegangen, unterwegs oder ausgeführt und die Antwort verloren sein. Der Sender kann die drei Fälle nicht unterscheiden --- und muss trotzdem handeln.}
\end{center}

\begin{technik}
\boxtitel{Die Größenordnungen}
\begin{center}
\renewcommand{\arraystretch}{1.45}
\begin{tabular}{@{}P{3.0cm}P{2.2cm}P{2.6cm}@{}}
\toprule
\textbf{Vorgang} & \textbf{Dauer} & \\
\midrule
Funktionsaufruf & $\approx2$ ns & \\
Sprachgrenze & $\approx60$ ns & \textbf{Band 21, Kap.~14}\\
Prozessgrenze & $\approx5$ µs & \textcolor{mittel}{Rohr, Steckdose}\\
Rechner im Haus & $\approx0{,}5$ ms & \textbf{Faktor $\mathbf{250\,000}$}\\
Rechenzentrum entfernt & $\approx80$ ms & \\
\bottomrule
\end{tabular}
\end{center}

\begin{center}
\begin{minipage}{0.86\textwidth}
\kursiv{Zwischen einem Funktionsaufruf und einem Aufruf über das Hausnetz liegt Faktor zweihundertfünfzigtausend. Wer eine Naht dort legt, wo vorher ein Aufruf stand, ändert nicht die Struktur, sondern die Größenordnung --- und Größenordnungen ändern die richtige Bauform.}
\end{minipage}
\end{center}
\end{technik}
\end{merksatz}

\begin{merksatz}
\boxtitel{Die Zerlegung kommt vor der Verteilung}
\begin{center}
\renewcommand{\arraystretch}{1.55}
\begin{tabular}{@{}P{2.8cm}P{6.0cm}@{}}
\toprule
\textbf{Zuerst} & Wonach schneidet man? --- Parnas' Frage aus Band 17\\
\textbf{Dann erst} & \textbf{Was davon läuft in einem eigenen Prozess?}\\
\textbf{Der Fehler} & Beides zugleich entscheiden\\
\bottomrule
\end{tabular}
\end{center}

\begin{center}
\begin{minipage}{0.88\textwidth}
\kursiv{Ein schlecht geschnittenes Programm wird durch Aufteilen zu einem schlecht geschnittenen Werk --- mit denselben Abhängigkeiten, nun über das Netz, mit Latenz und Teilausfall. Man hat nichts gewonnen und fünf Schwierigkeiten hinzugefügt.}
\end{minipage}
\end{center}

\textbf{Die Probe.}

\begin{center}
\renewcommand{\arraystretch}{1.5}
\begin{tabular}{@{}P{3.0cm}P{2.2cm}P{2.8cm}@{}}
\toprule
\textbf{Frage} & \textbf{Antwort} & \\
\midrule
Ließe sich der Teil als Modul sauber trennen? & \textcolor{gruen}{\textbf{muss ja sein}} & \textbf{sonst nicht verteilen}\\
Ändert er sich aus eigenem Grund? & \textcolor{gruen}{muss ja sein} & \\
Hat er eigene Betriebsanforderungen? & \textcolor{gruen}{dann lohnt es} & \\
Nur weil es modern ist? & \textcolor{akzent}{\textbf{nein}} & \\
\bottomrule
\end{tabular}
\end{center}

\begin{center}
\kursiv{Die dritte Zeile ist der beste positive Grund: Wenn ein Teil andere Zeitschranken, andere Verfügbarkeit oder anderen Speicherbedarf hat als der Rest, gehört er in einen eigenen Prozess. Beim Messdatenerfasser ist das offensichtlich; bei einer Auswertungsfunktion meist nicht.}
\end{center}

\begin{warnung}
\boxtitelr{Conways Beobachtung}
\begin{center}
\renewcommand{\arraystretch}{1.5}
\begin{tabular}{@{}P{2.8cm}P{6.0cm}@{}}
\toprule
\textbf{Die Aussage} & Ein System bildet die Kommunikationsstruktur seiner Erbauer ab\\
\textbf{Die Folge} & \textbf{Drei Teams erzeugen drei Bestandteile --- auch wenn zwei genügten}\\
\textbf{Der Umkehrschluss} & Wer die Struktur ändern will, ändert zuerst die Zuständigkeiten\\
\bottomrule
\end{tabular}
\end{center}

\begin{center}
\begin{minipage}{0.88\textwidth}
\kursiv{Die Beobachtung ist fast sechzig Jahre alt und in jedem größeren Vorhaben nachprüfbar. Sie erklärt, warum Nähte oft dort verlaufen, wo Abteilungsgrenzen liegen --- und nicht dort, wo die Sache es verlangt.}
\end{minipage}
\end{center}

\begin{center}
\kursiv{Für ein Klinikvorhaben heißt das konkret: Wenn Labor, Station und Verwaltung je eigene Zuständigkeiten haben, entstehen drei Bestandteile mit drei Nähten --- unabhängig davon, ob der Befund das nahelegt.}
\end{center}
\end{warnung}
\end{merksatz}

\begin{aufgabe}
\begin{enumerate}[label=(\alph*)]
\item Bewerten Sie eine vorgeschlagene Aufteilung.
\item Rechnen Sie die Kosten einer Naht aus.
\item Bestimmen Sie, was ein Werk fertig macht.
\end{enumerate}

\loesung
\textbf{(a) Eine Aufteilung bewerten.}

\textbf{Der Vorschlag.} Die Auswertung in vier Dienste zerlegen.

\begin{center}
{\small
\renewcommand{\arraystretch}{1.45}
\begin{tabular}{@{}P{2.8cm}P{2.0cm}P{2.0cm}P{1.8cm}@{}}
\toprule
\textbf{Vorschlag} & \textbf{Eigener Grund?} & \textbf{Eigener Betrieb?} & \textbf{Urteil}\\
\midrule
Befunddienst & \textcolor{gruen}{ja} & \textcolor{gruen}{ja} & \textcolor{gruen}{\textbf{ja}}\\
Verlaufsdienst & \textcolor{akzent}{\textbf{nein}} & \textcolor{akzent}{nein} & \textcolor{akzent}{\textbf{nein}}\\
Anlassdienst & \textcolor{akzent}{nein} & \textcolor{akzent}{nein} & \textcolor{akzent}{nein}\\
Benachrichtigung & \textcolor{gruen}{ja} & \textcolor{gruen}{\textbf{ja}} & \textcolor{gruen}{ja}\\
\bottomrule
\end{tabular}}
\end{center}

\begin{center}
\begin{minipage}{0.88\textwidth}
\kursiv{Zwei von vier tragen. Verlauf und Anlassbestimmung sind reine Rechnung auf denselben Daten und ändern sich aus demselben Grund --- sie gehören in denselben Prozess. Ein eigener Dienst dafür erzeugt zwei Nähte für null Gewinn.}
\end{minipage}
\end{center}

\textbf{Die Rechnung für den Verlaufsdienst.}

\begin{center}
\renewcommand{\arraystretch}{1.45}
\begin{tabular}{@{}P{2.8cm}P{2.2cm}P{2.4cm}P{1.8cm}@{}}
\toprule
& \textbf{Im Prozess} & \textbf{Als Dienst} & \\
\midrule
Aufrufdauer & $\mathbf{48}$ ms & $\mathbf{51}$ ms & \\
Bei Ausfall & \textcolor{gruen}{unmöglich} & \textcolor{akzent}{\textbf{Auswertung steht}} & \\
Zeilen zusätzlich & $0$ & $\approx340$ & \\
Bereitstellungen & $1$ & $2$ & \\
Fehlerklassen neu & $0$ & $\mathbf{5}$ & \\
\bottomrule
\end{tabular}
\end{center}

\begin{center}
\kursiv{Drei Millisekunden mehr, dreihundertvierzig Zeilen Netzcode, eine zusätzliche Ausfallmöglichkeit und fünf neue Fehlerklassen --- für einen Bestandteil, der mit dem Aufrufer immer gemeinsam geändert wird.}
\end{center}

\medskip
\textbf{(b) Kosten einer Naht.}

\begin{center}
{\small
\renewcommand{\arraystretch}{1.45}
\begin{tabular}{@{}P{3.0cm}P{2.4cm}P{2.6cm}@{}}
\toprule
\textbf{Posten} & \textbf{Einmalig} & \textbf{Laufend}\\
\midrule
Format vereinbaren & $\approx2$ PT & \\
Beide Seiten umsetzen & $\approx5$ PT & \\
Fehlerbehandlung & $\approx4$ PT & \\
Prüfung über die Naht & $\approx3$ PT & \textcolor{akzent}{\textbf{je Änderung}}\\
Betriebs\-überwachung & $\approx2$ PT & \textcolor{akzent}{Bereitschaft}\\
Versionierung & $\approx2$ PT & \textcolor{akzent}{\textbf{dauerhaft}}\\
\midrule
\textbf{Summe} & $\approx\mathbf{18}$ \textbf{PT} & \\
\bottomrule
\end{tabular}}
\end{center}

\begin{center}
\begin{minipage}{0.88\textwidth}
\kursiv{Achtzehn Personentage für eine Naht, bevor eine Zeile Fachlogik geschrieben ist. Bei fünf Bestandteilen und vier Nähten sind das zweiundsiebzig Tage --- ein Drittel Personenjahr, das im Zeitplan meist nicht auftaucht.}
\end{minipage}
\end{center}

\begin{center}
\kursiv{Die vorletzte Zeile ist die dauerhafteste: Eine Naht muss versioniert werden, weil beide Seiten nicht gleichzeitig umgestellt werden können. Kapitel~\ref{kap:vertraege} behandelt, wie.}
\end{center}

\medskip
\textbf{(c) Was ein Werk fertig macht.}

\begin{center}
{\small
\renewcommand{\arraystretch}{1.5}
\begin{tabular}{@{}P{3.2cm}P{2.2cm}P{2.6cm}@{}}
\toprule
\textbf{Bedingung} & \textbf{Bei einem Programm} & \textbf{Bei einem Werk}\\
\midrule
Rechnet richtig & \textcolor{gruen}{\textbf{ja}} & \textcolor{gruen}{ja}\\
Läuft schnell genug & \textcolor{gruen}{ja} & \textcolor{gruen}{ja}\\
Prüfbar & \textcolor{gruen}{ja} & \textcolor{gruen}{ja}\\
\textbf{Übersteht Teilausfall} & --- & \textcolor{gruen}{\textbf{ja}}\\
\textbf{Im Betrieb befragbar} & --- & \textcolor{gruen}{\textbf{ja}}\\
\textbf{Einzeln erneuerbar} & --- & \textcolor{gruen}{\textbf{ja}}\\
\bottomrule
\end{tabular}}
\end{center}

\begin{center}
\begin{minipage}{0.88\textwidth}
\kursiv{Drei zusätzliche Bedingungen gegenüber Band 20, Kapitel~14 --- und alle drei sind prüfbar. Die fünfte wird am häufigsten vergessen: Ein Werk, das man im Betrieb nicht befragen kann, ist nicht fertig, gleichgültig wie gut es rechnet.}
\end{minipage}
\end{center}

\begin{center}
\kursiv{Die sechste ist die anspruchsvollste: Jeder Bestandteil muss einzeln ausgetauscht werden können, während die anderen laufen. Wer das nicht kann, muss für jede Änderung das ganze Werk anhalten --- in einer Klinik ein Vorgang mit Terminplanung.}
\end{center}

\begin{warnung}
\boxtitelr{Der Zeitplan lügt bei Werken}
\begin{center}
\renewcommand{\arraystretch}{1.45}
\begin{tabular}{@{}P{2.8cm}P{2.2cm}P{2.4cm}P{1.8cm}@{}}
\toprule
\textbf{Posten} & \textbf{Geschätzt} & \textbf{Tatsächlich} & \\
\midrule
Fachlogik & $60\,\%$ & $\mathbf{35\,\%}$ & \\
Nähte & \textcolor{akzent}{$15\,\%$} & \textcolor{akzent}{$\mathbf{30\,\%}$} & \\
Betrieb, Überwachung & $10\,\%$ & $\mathbf{20\,\%}$ & \\
Fehlerfälle & $15\,\%$ & $15\,\%$ & \\
\bottomrule
\end{tabular}
\end{center}

\begin{center}
\begin{minipage}{0.88\textwidth}
\kursiv{Die zweite und dritte Zeile werden regelmäßig halbiert geschätzt, weil sie beim Entwurf unsichtbar sind: Man plant, was das System \kursiv{tun} soll, und nicht, wie die Teile miteinander reden und wie man ihnen dabei zusieht.}
\end{minipage}
\end{center}

\begin{center}
\kursiv{Brooks' Beobachtung, dass Vorhaben nicht durch mehr Leute schneller werden, hat hier eine Entsprechung: Sie werden auch nicht dadurch schneller, dass man sie in mehr Bestandteile zerlegt. Jede Naht ist Kommunikationsaufwand, in Code gegossen.}
\end{center}
\end{warnung}
\end{aufgabe}
\vspace{-4pt}
\noindent Parnas\zit{19} stellt die Zer\-le\-gungs\-fra\-ge, die vor jeder Ver\-tei\-lung steht; Ford und Richards\zit{9} ordnen die Ent\-wurfs\-ent\-schei\-dun\-gen im Großen; Brooks\zit{5} beschreibt, warum Auf\-tei\-lung Kom\-mu\-ni\-ka\-tions\-auf\-wand erzeugt.
